Fix a score vector outside the \(P_X^{\otimes n}\)-null set on which either a conditional kernel assertion or the conditional product factorization below can fail.

First consider the deterministic trace count.  For a closed disk with center \(c\in\mathbb R^2\) and radius \(r\), membership of a fixed point \(X_i\) is equivalent to
\[
 2c^{\mathsf T}X_i+s\ge \lVert X_i\rVert_2^2,
 \qquad s=r^2-\lVert c\rVert_2^2. \tag{10}
\]
Thus every disk trace on \(\{X_1,\ldots,X_n\}\) is a sign pattern induced by \(n\) affine hyperplanes in the three-dimensional parameter \((c_1,c_2,s)\).  A trace realized by a closed disk can be made strict without changing it: if some points are excluded, increase the radius by less than their strictly positive minimum distance from the disk; if none is excluded, use any sufficiently small increase.  Equality in (10) therefore creates no additional trace.

Let \(A(n,v)\) be the largest number of open cells cut out by \(n\) affine hyperplanes in \(\mathbb R^v\).  Adding the last hyperplane retains all old cells and splits at most \(A(n-1,v-1)\) of them, because the arrangement induced on that hyperplane has dimension \(v-1\).  Hence
\[
 A(n,v)\le A(n-1,v)+A(n-1,v-1),
 \qquad A(0,v)=A(n,0)=1. \tag{11}
\]
Induction using Pascal's identity gives
\[
 A(n,3)\le\sum_{j=0}^3\binom nj. \tag{12}
\]
Restricting to genuine disks, fixing \(h\), restricting centers to \(\mathcal B\), or intersecting each trace with the fixed side set \(\{i:X_i\in\mathcal A_t\}\) cannot increase the count.  Multiplication by the two choices of \(t\) and the finite cardinality \(|\mathcal H_n|\), supplied by \(\mathrm{def:dyadic\text{-}grid}\), proves
\[
 |\mathscr S_n|\le2|\mathcal H_n|\sum_{j=0}^3\binom nj. \tag{13}
\]

If \(\mathscr S_n=\varnothing\), the stochastic conclusions follow from the stipulated zero-maximum convention.  Suppose it is nonempty.  By \(\mathrm{ass:iid\text{-}sampling}\), the observed pairs are independent, and regular conditional laws given \((X_1,\ldots,X_n)\) factor for \(P_X^{\otimes n}\)-almost every score vector.  Hence the residuals \(\varepsilon_i=Y_i-g_{T_i,P}(X_i)\) are conditionally independent.  Because \(\mathcal A_0=\mathbb R^2\setminus\mathcal A_1\), an index in a side-\(t\) candidate set satisfies \(T_i=t\).  The restricted score measure assigns full mass to \(\mathcal X_{t,P}\): its complement is a countable union of rational basic open sets having restricted mass zero.  Taking the union of the two fixed null sets in \(\mathrm{ass:gaussian\text{-}errors}\), that assumption applies to every relevant index and gives, for all \(\lambda\in\mathbb R\),
\[
 \mathbb E_P\{e^{\lambda\varepsilon_i}\mid X_1,\ldots,X_n\}
 =\int e^{\lambda u}K_{T_i,P}(X_i,du)
 \le e^{\sigma^2\lambda^2/2}. \tag{14}
\]
The same assumption supplies conditional integrability and centering; no Gaussian equality is used.

For nonempty \(S\in\mathscr S_n\), set
\[
 Z_S=\frac1{\sigma\sqrt{|S|}}\sum_{i\in S}\varepsilon_i.
\]
Conditional independence and (14) imply, for every \(u\in\mathbb R\),
\[
 \mathbb E_P\{e^{uZ_S}\mid X_1,\ldots,X_n\}
 \le\prod_{i\in S}e^{u^2/(2|S|)}=e^{u^2/2}. \tag{15}
\]
For \(z>0\), Markov's inequality applied to \(e^{zZ_S}\), followed by (15), gives
\[
 \Pr_P\{Z_S>z\mid X_1,\ldots,X_n\}\le e^{-z^2/2}.
\]
The same calculation for \(-Z_S\), and then a union bound over \(S\), yields
\[
 \Pr_P\left\{\max_{S\in\mathscr S_n}|Z_S|>z
 \ \middle|\ X_1,\ldots,X_n\right\}
 \le2|\mathscr S_n|e^{-z^2/2}. \tag{16}
\]
No independence among the set-indexed sums is used.

Finally, let \(z_0=\sqrt{2\log\{2|\mathscr S_n|\}}\).  Since \(|\mathscr S_n|\ge1\), one has \(z_0>1\).  The tail-integral identity for a nonnegative random variable, (16), and
\[
 \int_{z_0}^{\infty}e^{-z^2/2}\,dz\le z_0^{-1}e^{-z_0^2/2}
\]
give
\[
 \begin{aligned}
 \mathbb E_P\left[\max_{S\in\mathscr S_n}|Z_S|
 \ \middle|\ X_1,\ldots,X_n\right]
 &\le z_0+2|\mathscr S_n|\int_{z_0}^{\infty}e^{-z^2/2}\,dz\\
 &\le z_0+z_0^{-1}\le z_0+1.
 \end{aligned}
\]
Together with the empty-family case, this is the claimed expectation bound with \(1\vee|\mathscr S_n|\).